\section{Application: Low density parity check or expander codes}
\label{sec:codes}

Another useful application of random graphs is as expanders,
loosely defined as graphs where the neighborhood of a small set of nodes is large.
Expander codes, also called low density parity check (LDPC) codes,
were first introduced by Gallager in his PhD thesis
(\citet*{gallager1962}).
These are a family of linear error correcting codes whose parity-check matrix 
is encoded in an expander graph.
A linear code is a set $\mathcal{C} \subset \Sigma^L$,
where a length $L$ codeword
$x \in \mathcal{C}$ if and only if
$H x = 0$.
The alphabet $\Sigma$ is typically a finite field
and $H \in \Sigma^{P \times L}$ is the parity check matrix.
In the simplest case, $\Sigma = \mathbb{F}_2$ and each row of $H$ 
can be interpreted as a parity constraint on codewords.
The performance of such codes depends on how good an expander that graph is,
which in turn can be shown to depend on the separation of eigenvalues.
For a good introduction and overview of the subject, 
see \citet*{richardson2008}.

Following \citet*{tanner1981}, 
we construct a code $\mathcal{C}$ 
from a $(d_1,d_2)$-regular bipartite graph $G$
on $n + m$ vertices
and two smaller linear codes 
$\mathcal{C}_1$ and $\mathcal{C}_2$ of length
$d_1$ and $d_2$, respectively.
We write 
$\mathcal{C}_1 = [d_1, k_1, \delta_1]$ 
and
$\mathcal{C}_2 = [d_2, k_2, \delta_2]$
with the usual convention
of length, dimension, and minimum distance.
We assume the codes are all binary, using the finite field $\mathbb{F}_2$ 
the codeword is $x \in \mathcal{C} \subset \mathbb{F}_2^{|E|}$
where  $|E| = n d_1 = m d_2$.
That is, we associate a bit to each edge in the graph bipartite graph $G$.
Let $( e_i(v) )_{i=1}^{d_v}$ 
represent the set of edges incident to a vertex $v$
in some arbitrary, fixed order.
Then the vector $x \in \mathcal{C}$ if and only if the vectors
$( x_{e_1(u)}, x_{e_2(u)}, \ldots, x_{e_{d_1}(u)} )^T \in \mathcal{C}_1$
for all $u \in V_1$ and
$( x_{e_1(v)}, x_{e_2(v)}, \ldots, x_{e_{d_2}(v)} )^T \in \mathcal{C}_2$
for all $v \in V_2$.
The final code $\mathcal{C}$ is also linear.
With this construction, the code $\mathcal{C}$ 
has rate at least $k_1 / d_1 + k_2/d_2 - 1$ (\citet*{tanner1981}).

Furthermore, \citet*{janwa2003} proved the following bound on
the minimum distance of the resulting code:
\begin{theorem}
\label{thm:JanwaLal}
Suppose $\delta_1 \geq \delta_2 > \eta/2$, where $\eta$ is the second largest eigenvalue of the adjacency matrix of $G$. Then the code $\mathcal{C}$ has minimum distance
\[
\delta \geq \frac{n}{d_2} \left( \delta_1 \delta_2 - \frac{\eta}{2} (\delta_1 + \delta_2 ) \right).
\]

\end{theorem}

\begin{corollary}
\label{cor:dist_bound}
Suppose the code $\mathcal{C}$ is constructed from a biregular, bipartite random graph 
$G \sim \mathcal{G}(n,m, d_1, d_2)$ and the conditions of Theorem~\ref{thm:JanwaLal} hold.
Then the minimum distance of $\mathcal{C}$ satisfies
\[
\delta \geq \frac{n}{d_2} \left( \delta_1 \delta_2 - 
\frac{\sqrt{d_1 - 1} + \sqrt{d_2 - 1}}{2} (\delta_1 + \delta_2 ) - \epsilon_n \right) .
\]
\end{corollary}

We see that these Tanner codes will have maximal distance for smallest $\eta$,
and used our main result, Theorem~\ref{cor:gap_A}, to obtain the explicit
bound in Corollary~\ref{cor:dist_bound}.
By growing the graph, the above shows a way to construct arbitrarily large codes whose
minimum distance remains proportional to the code size $n d_1$.
That is, the relative distance $\delta / (n d_1)$ 
is bounded away from zero as $n \to \infty$.
However, the above bound will only be useful if it yields a positive result, 
which depends on the codes $\mathcal{C}_1$ and $\mathcal{C}_2$ as well as the degrees.

\begin{remark}
In general, the performance guarantees on LDPC codes that are obtainable 
from graph eigenvalues are weaker than those that come from other methods.
Although our method does guarantee high distance for some high degree codes, 
analysis of specific decoding algorithms or a probabilistic expander analyses
yield better bounds that work for lower degrees (\citet*{richardson2008}).
\end{remark}

\subsection{Example: An unbalanced code based on a $(14,9)$-regular bipartite graph}

We illustrate the applicability of our distance bound with an example.
Let $\mathcal{C}_1 = [14, 8, 7]$ and $\mathcal{C}_2 = [9, 4, 6]$.
These can be achieved by using a Reed-Salomon code
on the common field $\mathbb{F}_{q}$ for any $q > 14$ (\citet*{richardson2008}).
We take $q = 2^4 = 16$ for inputs that are actually binary, and
this means each edge in the graph actually contains 4 bits of information.
Employing Corollary~\ref{cor:dist_bound}, 
the Tanner code $\mathcal{C}$
will have relative minimum distance 
$\delta / (n d_1) \geq 0.0014$ and 
rate at least $0.016$.
Taking $n = 216$ and $m = 336$ gives the code a minimum distance of 
at least 4.











